\chapter{Хэллэг ба Логик, Үнэний хүснэгт}
\section*{Лекцийн зорилго}
\begin{itemize}
  \item Логик, хэллэгийн үндсэн ойлголтыг тодорхойлох.
  \item Үнэний хүснэгтээр логик томьёог шинжлэн үнэнийг тогтоох аргыг эзэмшүүлэх.
  \item Баталгааны үндсэн аргуудыг тодорхойлолт бүрээр нь жишээтэй тайлбарлах.
  \item Олонлог, түүн дээрх үйлдлүүд
\end{itemize}
Цаашид хэрэглэгдэх тэмдэглэгээнүүд\\
\begin{multicols}{2}
\begin{enumerate}
\item $\forall$- дурын, аливаа
\item $\exists$- оршин байна, олдоно
\item $\exists!$- цор ганц оршин байна, ганц олдоно 
\item $\Rightarrow$- мөрдөнө, мөрдөн гарна
\item $\in$ - харяалагдна
\item $\notin$- харяалагдахгүй
\item $\Leftrightarrow$- эквивалент, зайлшгүй бөгөөд хүрэлцээтэй
\item $\mathbb{N}$- натурал тоон олонлог
\item $\mathbb{Z}$- бүхэл тоон олонлог
\item $\mathbb{Q}$-рационал тоон олонлог
\item $\mathbb{R}$- бодит тоон олонлог
\end{enumerate}
 Грек үсгүүдийн нэршил\\
\begin{tabular}{c|c|l}
\hline
Том үсэг & Жижиг үсэг & Дуудлага \\ \hline
$A$ & $\alpha$ & альфа \\
$B$ & $\beta$ & бета \\
$\Gamma$ & $\gamma$ & гамма \\
$\Delta$ & $\delta$ & дельта \\
$E$ & $\epsilon$ & эпсилон \\
$Z$ & $\zeta$ & зета \\
$H$ & $\eta$ & эта \\
$\Theta$ & $\theta$ & тета \\
$\Xi$ & $\xi$ & кси \\
$\Lambda$ & $\lambda$ & лямбда \\
$M$ & $\mu$ & мю \\
$N$ & $\nu$ & ню \\
$\Sigma$ & $\sigma$ & сигма \\
$T$ & $\tau$ & тау \\
$\Phi$ & $\varphi$ & фи \\
$\Psi$ & $\psi$ & пси  \\
\hline
\end{tabular}
\end{multicols}
\begin{remark}
     $\forall$, $\exists$, $\exists!$, $\Rightarrow$, $\in$, $\Leftrightarrow$ тэмдэгүүдийг квантор гэж нэрлэдэг.
\end{remark}
\begin{example} Логик хэллэгийн уншилтын жишээ
  \begin{enumerate} 
      \item
\[
\forall x\in\mathbb{R} \text{ бол } \big( x>1  \Rightarrow  x^2>1 \big).
\]
\textit{Уншилт:} Бүх бодит \(x\)-ийн хувьд, хэрэв \(x>1\) бол \(x^2>1\).\par\medskip
      \item \[
\exists!\,n\in\mathbb{N}\, , \,\big( n^2=25\big).
\]
\textit{Уншилт:} Натурал тоонууд дундаас \(n^2=25\) байх \emph{цор ганц} \(n\) (өөрөөр хэлбэл \(n=5\)) оршин байна.\par\medskip
      \item\[
\forall n\in\mathbb{Z}\ \text{ хувьд } \big(n \text{ тэгш}\big)\ \Leftrightarrow\ \big(\exists k\in\mathbb{Z}\ n=2k\big).
\]
\textit{Уншилт:} Бүх бүхэл \(n\)-ийн хувьд, \(n\) тэгш байх \emph{зайлшгүй бөгөөд хүрэлцээтэй нөхцөл} нь зарим бүхэл \(k\) тооны хувьд \(n=2k\) гэж бичигдэх нөхцөл болно.\par\medskip
\textit{Өөрөөр хэлбэл :} Бүх бүхэл \(n\)-ийн хувьд, \(n\)  \emph{зөвхөн тэгш байх үед л} зарим бүхэл \(k\) тооны хувьд \(n=2k\) гэж бичигдэх болно.\par\medskip
      \item \[
\forall \epsilon>0\, ,  \exists \delta>0\,  , \forall x\in\mathbb{R}\ 
\Big(|x-1|<\delta\ \Rightarrow\ \big|\tfrac{x^2-1}{x+1}-0\big|<\epsilon\Big).
\]
\textit{Уншилт:} Дурын эерэг \(\epsilon\) тооны хувьд, ямар нэг эерэг \(\delta\) тоо олдох ба түүний хувьд, бүх бодит \(x\)-ийн хувьд \(|x-1|<\delta\) бол
\(\big|\frac{x^2-1}{x+1}\big|<\epsilon\) байна.\par\medskip
  \end{enumerate}  
\end{example}
\newpage
\section{Хэллэг ба логик}
\subsection{Хэллэг}
\begin{definition} \emph{Үнэн эсвэл худал гэдгийг нь нэгэн утгатай тодорхой тогтоож болох} өгүүлбэрийг \textbf{хэллэг} (proposition) гэнэ.
\end{definition}
 Өөрөөр хэлбэл, тухайн хэллэгийн үнэний утга (үнэн/худал) нь хамрах хүрээ, хугацаа, хэрэглэсэн ухагдахуунууд тодорхой үед ганц утгатайгаар тодорхойлогддог хэллэг юм.
 \begin{example}
     \begin{itemize}
  \item ``$5<9$'' — үнэн хэллэг.
  \item ``7 нь тэгш.'' — худал хэллэг.
  \item ``$\pi$ нь рационал тоо.'' — худал хэллэг.
  \item ``бүх $x\in\mathbb{R}$ хувьд $x^2\ge 0$ .'' — үнэн хэллэг (квантортой хэллэг).
  \item ``2026 оны 1-р сарын 1-нд МУИС I байрны цэцэрлэгт цас орсон.'' — хугацаа, газар тодорхой тул үнэн эсэхийг тогтоож болно.
\end{itemize}
 \end{example}
\begin{remark}
(1) Тушаал, хүсэлт, асуулт гэх мэт нь хэллэг биш: \\
``Хаалгаа хаа!'', \\
``Хаалга хаасан уу?''. \\
(2) \emph{Нээлттэй хэллэг} буюу хувьсагчтай, кванторгүй хэллэг (ж: ``$x>2$'') нь хэллэг биш; \\
харин $\forall x\in\mathbb{R}\,(x>2 \Rightarrow x^2>4)$ шиг квантортой бол хэллэг болно. \\
(3) Дунджаар, ``дулаахан'' гэх мэт тодорхойгүй шалгууртай, контекст тодорхойгүй хэллэгүүд нь хэллэг биш; шалгуур, хэмжүүрийг тодорхой тогтоосноор хэллэг болгож болно.
\end{remark}

\begin{definition}
Хэллэг үнэн бол түүнийг 1, худал бол түүнийг 0 утгатай гээд 0 ба 1 утгийг хэллэгийн \textbf{үнэний утга} гэнэ.
\end{definition}
Бид цаашид хэллэгийг үсгээр нэрлэх болно. 
$$P=\{11>9\} \,,\, Q=\{\exists n\in\mathbf{N} ,  \big( n^2=-1\big)\}$$ гэе. Тэгвэл $P$ хэллэг үнэн тул 1, $Q$ хэллэг худал тул 0 утгатай.
\subsection{Логик холбоосууд}
Бид олон хэллэгүүдийг нийлүүлэн тэдгээрээс дүгнэлт гаргадаг. Өөрөөр хэлбэл\\ математикийн теорем гэдэг нь нийлмэл хэллэг байдаг. Бид энэ хэсэгт нийлмэл хэллэг хэзээ үнэн байх тухай авч үзнэ. 

Эхлээд нэг эсвэл хэд хэдэн хэллэгээс хэрхэн шинэ хэллэг үүсгэх талаар авч үзье.

\begin{definition}[Үгүйсгэл] Зөвхөн $A$ хэллэг худал байхад үнэн байдаг хэллэгийг $\lnot A$ гэж тэмдэглэх ба "$A$ биш" гэсэн хэллэг болно. Үүнийг үгүйсгэл гэнэ. 
\end{definition}
\begin{definition}[ Импликаци, “Хэрэв…бол”] Зөвхөн $A$ хэллэг худал байхад эсвэл $B$ хэллэг үнэн байхад үнэн байдаг хэллэгийг $A\Rightarrow B$ гэж тэмдэглэх ба  "$A$ байсан бол $B$  байна" гэсэн хэллэг болно. $B$ хэллэгийг $A$ хэллэгийн мөрдөлгөө гэнэ. Энэ нийлмэл хэллэгээр шалтгаан үр дагаварын холбоог илэрхийлнэ.
\end{definition}
\begin{definition}[Үржвэр, "... ба ..."]
$A$ ба $B$ хэллэгүүд хоёул үнэн байхад үнэн байдаг бусад үед худал байдаг  хэллэгийг $A {\text{ ба }} B$ эсвэл $A\land B$ гэж тэмбэглэх ба энэ нь "$A$ ба $B$" гэсэн хэллэг болох ба $A$, $B$ хэллэгүүдийн үржвэр /conjuction / гэдэг. 
\end{definition}
\begin{definition}[Нийлбэр, "... эсвэл ..."]
$A$ ба $B$ хэллэгүүдийн ядаж нэг үнэн байхад үнэн байдаг бусад үед худал байдаг хэллэгийг $A\lor B$ гэж тэмдэглэх ба энэ нь $"A {\text{ эсвэл }} B" $  гэсэн хэллэг болох ба $A$, $B$ хэллэгүүдийн  нийлбэр /disjunction / гэдэг. 
\end{definition}
Мөн нь нийлмэл хэллэгийн үнэн-утгыг хувьсагчдын (хэллэгүүдийн) бүх боломжит үнэн худал утгийн хослол бүр дээр системтэйгээр тооцож харуулсан хүснэгтийг нийлмэл хэллэгийн үнэний хүснэгт гэдэг.


$\lnot A$ , $A\lor B$, $A\land B$ , $A\Rightarrow B$ үзэгдлүүдийн үнэний хүснэгтийг хийж үзье. \\
$$\begin{tabular}{|c|c|c|c|c|c|}
\hline
A&B&$\lnot$ A&A\lor B&A\land B&$A\Rightarrow B$\\
\hline
1&1&0&1&1&1\\
\hline
1&0&0&1&0&0\\
\hline
0&1&1&1&0&1\\
\hline
0&0&1&0&0&1\\
\hline
\end{tabular}$$
\begin{example}
$$A=\{{\text{Бороо орсон}}\}$$
$$B=\{{\text{Газар нойтон байна}}\}$$ гэсэн хэллэгүүдийг авч үзье. 
\begin{itemize}
\item $\lnot A=\{{\text{Бороо ороогүй}}\}$\\
Энэ хэллэг $A$ тай яг эсрэг байна.
\item $A\lor B=\{{\text{Бороо орсон эсвэл Газар нойтон байна}}\}$\\
Энэ нийлмэл хэллэгт орсон 2 хэллэгийн ядаж 1 нь үнэн үед $A\lor B$-г үнэн гэж үзнэ. Өөрөөр хэлбэл "Бороо орсон боловч газар хуурай байсан ч", "Бороо ороогүйч газар нойтон байсан ч", "Бороо ороод газар нойтон байсан ч" "$A$ эсвэл $B$" хэллэг үнэн байна. 
\begin{remark}
    Нөхцөл $A$ худал үед $A\Rightarrow B$ үргэлж үнэн гэж үздэг. 
    
    Жишээ: ``Хэрэв би өнөөдөр ирвэл (A) амттан авна (B).'' Би ирээгүй (A худал) бол амттан авсан/аваагүйгээс үл хамааран хэллэгийг худал болгох шалтгаан байхгүй тул нийт хэллэг үнэн гэж тооцно.
\end{remark}
\item $A\land B=\{{\text{Бороо орсон ба Газар нойтон байна}}\}$\\
Энэ нийлмэл хэллэгт орсон 2 хэллэг 2-уул үнэн үед $A\land B$-г үнэн гэж үзнэ. Өөрөөр хэлбэл "Бороо орсон бас газар нойтон" байсан үед л үнэн бусад үед худал гэж үзнэ.
\item $(A\Rightarrow B)=\{{\text{Бороо орсон бол Газар нойтон байна}}\}$\\
Энэ нийлмэл хэллэгт Бороо орсон нь шалтгаан,   Газар нойтон байх нь үр дүн юм. Энэ хэллэг нь "Бороо орсон бас   Газар нойтон биш" байхад л худал, бусад үед үнэн гэж үзнэ. 
\end{itemize}
\end{example}
\begin{definition}[Эквивалент, if and only if] 
$A\Rightarrow B$ мөн $B\Rightarrow A$ бол $A$ ба $B$-г эквивалент  /ижил/ хэллэгүүд ба $A$-г $B$-ийн зайлашгүй бөгөөд хүрэлцээтэй нөхцөл гээд $A\Leftrightarrow B$ гэж тэмдэглэнэ. 
\end{definition} 
$A\Leftrightarrow B$ бол $A$ болон $B$ хэллэгүүд ижил үнэний утгатай мөн ижил ижил үнэний утгатай бол $A\Leftrightarrow B$ байна. 
\begin{example}
$x=2\Rightarrow x^2=4$ хэллэг үнэн.\\
$x^2=4\Rightarrow x=-2$ хэллэг худал.\\
$x+5=y+3\Leftrightarrow x+2=y$ хэллэг үнэн. \\
$\{n>10$ ба $n<12 {\text{ бүхэл тоо}}\}\Leftrightarrow \{n=11\}$\\
$\{n\leq 10$ эсвэл $n\geq 12 {\text{ бүхэл тоо}}\}\Leftrightarrow \{n\neq 11\}$\\
\end{example}
\subsection{Логикийн хуулиуд, тавтологи ба зөрчил}
\begin{definition}[Тавтологи] Ямагт үнэн байдаг хэллэгийг тавтологи гэнэ.
\end{definition}
\begin{definition}[Зөрчил]
    Ямагт худал байдаг хэллэгийг зөрчил гэнэ.
\end{definition}
\begin{definition}[Баталгаа]
Өгөгдсөн суурь өгүүлбэрүүдээс (аксиом, тодорхойлолт, өмнө батлагдсан хэллэгүүдээс) \emph{логикийн дүгнэлтийн дүрмүүд}-ийг ($\land$, $\lor$, $\Rightarrow$, $\lnot$, $\Leftrightarrow$ гэх мэт) хэрэглэн зорьсон дүгнэлтэд хүрэх төгсгөлөг алхамуудын дараалалыг баталгаа гэнэ. Алхам бүр нь 

(i) суурь өгүүлбэр, эсвэл 
(ii) өмнөх мөрүүдээс дүрмээр гарсан шинэ өгүүлбэр байна.
\end{definition}
\begin{definition}[Tеорем]
Аксиом, тодорхойлолт, өмнө батлагдсан хэллэгүүдээс логикийн зөв дүгнэлтийн дүрмүүдийг ашиглан төгсгөлөг \emph{алхамчилсан баталгаагаар} гаргаж авсан үнэн хэллэгийг теорем гэнэ.
\end{definition}
 Теорем нь ерөнхийдөө “Хэрэв (нөхцөлүүд) бол (дүгнэлт)” хэлбэртэй байна.

\begin{example}
   Теорем.  Хэрэв $P$ хэллэг бол $Q=P\lor\lnot P$ хэллэг тавтологи байна.

   Баталгаа.  $Q=P\lor\lnot P$ хэллэгийн үнэний хүснэгтийг байгуулж үзье. 
   \[
\begin{array}{c|c|c}
P & \lnot P & P\lor\lnot P\\\hline
1 & 0 & 1\\
0 & 1 & 1
\end{array}
\]
Сүүлийн баганын нүд бүрт 1 гарч байгаа тул $Q$ нь $P$ хэллэгийн бүх боломжит утганд үнэн, өөрөөр хэлбэл тавтологи байна.
\end{example}
\begin{example}\label{Урвуу хэллэг}
 Теорем. $P\Rightarrow Q$ хэллэг ба $\lnot Q\Rightarrow \lnot P$ хэллэг эквивалент байна.

   Баталгаа.  $P\Rightarrow Q$ хэллэг ба $\lnot Q\Rightarrow \lnot P$ хэллэгийн үнэний хүснэгтийг байгуулж үзье. 
\begin{table}[h]
\centering
\begin{tabular}{c|c|c|c|c|c}
$P$ & $Q$ & $\neg P$ & $\neg Q$ & $P \Rightarrow Q$ & $\neg Q \Rightarrow \neg P$ \\ \hline
1   & 1   & 0        & 0        & \textbf{1}        & \textbf{1}        \\ 
1   & 0   & 0        & 1        & \textbf{0}        & \textbf{0}        \\ 
0   & 1   & 1        & 0        & \textbf{1}        & \textbf{1}        \\
0   & 0   & 1        & 1        & \textbf{1}        & \textbf{1}        \\ 
\end{tabular}
\caption*{$P \Rightarrow Q \Leftrightarrow \neg Q \Rightarrow \neg P$ баталгаа}
\end{table}
\end{example}
\newpage
\begin{theorem}[Хэллэгийн холбоосын чанарууд] \label{Хэллэгийн холбоосын чанарууд} $P, Q, R$ дурын хэллэгүүд, $T$ нь тавтологи, $F$ зөрчил бол дараах чанарууд биелнэ.
  \begin{enumerate}
      \item \textbf{Коммутатив чанар}: $P\land Q\Leftrightarrow Q\land P$, $P\lor Q\Leftrightarrow Q\lor P$
      \item \textbf{Ассоциатив чанар}: $(P\lor Q)\lor R\Leftrightarrow P\lor(Q\lor R)$, $(P\land Q)\land R\Leftrightarrow P\land(Q\land R)$
      \item \textbf{Дистрибутив чанар}: $P\land(Q\lor R)\Leftrightarrow (P\land Q)\lor(P\land R)$, $P\lor(Q\land R)\Leftrightarrow (P\lor Q)\land(P\lor R)$
      \item \textbf{Идентити чанар}: $P\lor F\Leftrightarrow P$, $P\land T\Leftrightarrow P$
      \item \textbf{Доминант чанар}: $P\lor T\Leftrightarrow T$, $P\land F\Leftrightarrow F$
      \item \textbf{Идемпотент чанар}: $P\lor P\Leftrightarrow P$, $P\land P\Leftrightarrow P$
      \item \textbf{Давхар үгүйсгэл}: $\lnot\lnot P\Leftrightarrow P$
      \item \textbf{Де Морганы хууль}: $\lnot(P\land Q)\Leftrightarrow(\lnot P\lor\lnot Q)$, $\lnot(P\lor Q)\Leftrightarrow(\lnot P\land\lnot Q)$
  \end{enumerate}  
\end{theorem}
\begin{prof}
Жишээ болгон $P\land Q\Leftrightarrow Q\land P$ теоремийг батлая. 
    \[
\begin{array}{c|c|c|c|c}
P & Q & P\land Q & Q\land P & (P\land Q)\Leftrightarrow(Q\land P)\\\hline
1 & 1 & 1 & 1 & 1\\
1 & 0 & 0 & 0 & 1\\
0 & 1 & 0 & 0 & 1\\
0 & 0 & 0 & 0 & 1
\end{array}
\]
3 ба 4-р баганаас харахад $P\land Q$ ба $Q\land P$ хэллэгүүдийн үнэний хүснэгт адил байна. Иймд тодорхойлолт ёсоор $P\land Q\Leftrightarrow Q\land P$ байна. Мөн сүүлийн багана бүх утга 1 тул \((P\land Q)\Leftrightarrow(Q\land P)\) нь тавтологи байна. 
\end{prof}
\begin{theorem}[Нийлмэл хэллэгийн үгүйсгэл]\label{Нийлмэл хэллэгийн үгүйсгэл}
Хэрэв $P, Q$ дурын хэллэгүүд бол 
$$\lnot(P\Rightarrow Q)\Leftrightarrow P\land\lnot Q$$ ба $$\lnot(P\Leftrightarrow Q)\Leftrightarrow(P\land\lnot Q)\lor(\lnot P\land Q)$$ байна.
    
\end{theorem}
\begin{prof}
    Баталгаа илэрхий.
\end{prof} 
\begin{example} Хэрэв Теорем \ref{Нийлмэл хэллэгийн үгүйсгэл}-д $P=\text{"би хичээлдээ очно"}$, $Q=\text{"багш оноо өгнө"}$ гэвэл $\lnot(P\Rightarrow Q)=\text{"Би хичээлдээ ирсэн боловч багш оноо өгөөгүй."}$ болно. 
\end{example}
\section{Баталгааны аргууд}
Аливаа теорем батлахад өгөгдсөн нөхцөл, өмнө баталсан үнэн үр дүн, (теорем) ашиглан эдгээр хэллэгүүд дээр логик дүгнэлт хувиргалт хийж батлах үр дүнг үнэн болохыг баталдаг. Баталгааг хийх хэд хэдэн арга байдаг бөгөөд өргөн хэрэглэддэг аргуудыг товч тайлбарлая.
\begin{enumerate}
    \item \it{Шууд баталах арга}
    
    Үр дүнг шууд логик үйлдэл, тодорхойлолт, өгөгдсөн нөхцөл, өмнө баталсан үнэн үр дүн ашиглан шууд батлах арга.
\begin{example}
Бүхэл тоог 2 оор үржүүлэхэд гарах тоог тэгш тоо, тэгш биш бүхэл тоог сондгой тоо гэдэг. Хэрэв $n$ тэгш бол $n^2$ тэгш байхыг батал. 
\begin{solution}
$n$ тоог $n=2k$ гэж бичиж болно. Иймд $n^2=4k^2=2(2k^2)$ нь тэгш тоо болж батлагдав.Учир нь $2k^2$ нь бүхэл тоо юм.
\end{solution}
\end{example}
\item \it{Урвуу батлах арга} 

Жишээ \ref{Урвуу хэллэг} ёсоор $P\Rightarrow Q$ ба  $\lnot Q\Rightarrow \lnot P$ эквивалент байдаг. Иймд $P\Rightarrow Q$ батлахын оронд $\lnot Q\Rightarrow \lnot P$ батлах арга. Зарим тохиолдолд $\lnot Q\Rightarrow \lnot P$ хэлбэрийг батлах нь илүү хялбар байдаг. $\lnot Q\Rightarrow \lnot P$ хэллэгийг $P\Rightarrow Q$ хэллэгийн урвуу хэллэг гэнэ.
\begin{example}
"Хэрэв $n^2$ тэгш бол $n$ тэгш." 
\begin{solution} $P=n^2$ тэгш, $Q=n$ тэгш гэсэн хэллэгүүд авч үзье. Бид $$P\Rightarrow Q$$ хэллэгийг үнэн болохыг батална. $\lnot P=n^2$ сондгой, $\lnot Q=n$ сондгой гэсэн хэллэгүүд болно. Иймд $P\Rightarrow Q$ батлахын оронд $\lnot Q\Rightarrow \lnot P$ батлах нь хялбар байна.
$n$ сондгой $\Rightarrow n=2k+1\Rightarrow n^2=4k^2+4k+1=2(2k^2+2k)+1$ сондгой. 

Бид $n$ сондгой бол  $n^2$ сондгой байхыг баталлаа. Иймд хэрэв $n^2$ тэгш бол $n$ сондгой байж болохгүй. Тиймээс эхний $P\Rightarrow Q$ хэллэг үнэн.
\end{solution}
\end{example}
\item \it{Эсрэгээс нь батлах арга} 

Теорем \ref{Нийлмэл хэллэгийн үгүйсгэл} ёсоор $$\lnot(P\Rightarrow Q)\Leftrightarrow P\land\lnot Q$$ тул батлахыг хүссэн $P\Rightarrow Q$ хэллэгийн оронд $P\land\lnot Q$ хэллэгийг үнэн гэж үзээд, түүнээс логик дүгнэлтүүд гарган  зөрчилд ($R\land\lnot R$ маягийн) эсвэл боломжгүй тэгшитгэлд (жишээ нь $0=1$) хүрч батлах арга юм. Зөрчил гарсан тул $P\land\lnot Q\Rightarrow \lnot(P\Rightarrow Q)$ худал,  тэгэхээр $P\Rightarrow Q$ үнэн. Өөрөөр хэлбэл үнэн хэллэгээс худал хэллэг гарахгүй тул $\lnot(P\Rightarrow Q)$ худал гэсэн үг.
\begin{example} $\sqrt{2}$ иррационал тоо юм.
\begin{solution}
    $P$: "$\sqrt{2}$ бол иррационал" гэе. Эсрэгээр нь таамаглая: $\sqrt{2}$ рационал буюу $\lnot P=\{\sqrt{2}=\dfrac{a}{b}\}$ ($a,b\in\mathbb{Z}$, $\text{ХИЕХ}(a,b)=1$) үнэн гэе. Тэгвэл $2=\dfrac{a^2}{b^2}$, өөрөөр хэлбэл $a^2=2b^2$. Эндээс $a^2$ тэгш тул $a$ тэгш, $a=2k$. Иймд $2b^2=4k^2 \Rightarrow b^2=2k^2$, иймээс $b$ бас тэгш. Ингэснээр $a$ ба $b$ хоёул тэгш гарч, $\text{ХИЕХ}(a,b)=1$ гэсэн нөхцөлтэй зөрчилдөнө. Иймд $\lnot P$  таамаг худал, тэгэхээр анхны таамаг $\sqrt{2}$ иррационал тоо болов.
\end{solution}\end{example}
\item \it{Эсрэг жишээ гаргах арга}  

"$P(x)$ хэллэг бүх $x$ хувьсагчийн хувьд үнэн" гэсэн хэллэгийг худал болохыг батлахдаа $P(x)$ худал байх тодорхой нэг $x$ хувьсагчийн утгийг олж заах арга юм.
\begin{example}
 "Бүх тэгш тоо 4-т хуваагдана."   
\end{example}
\begin{solution}
 6 нь тэгш боловч 4-т хуваагдахгүй. Иймд "Бүх тэгш тоо 4-т хуваагдана." хэллэг худал юм. 
\end{solution}
\item \it{Математик индукц}


$n\in\N$ ба $P(n)$ өгүүлбэр өгөгдөв. Хэрэв 
\begin{enumerate}
    \item $P(1)$ үнэн
    \item $\forall k\in\mathbb{N}$ хувьд $ P(k)$ үнэн бол $P(k+1)$ үнэн байдаг
\end{enumerate} чанартай байвал $\forall n\in\N$ хувьд $P(n)$ өгүүлбэр үнэн байна. 
\begin{example}
Дурын натурал $n$ тооны хувьд 
\begin{equation}
1 + 2 + 3 + \dots + n = \frac{n(n+1)}{2}
\end{equation}
адилтгал биелнэ.
\begin{prof}
Математик индукцийн аргыг ашиглая.

\begin{enumerate}
    \item \textbf{Индукцийн суурь ($n=1$):} 
    Зүүн гар тал: $1$ \\
    Баруун гар тал: $\dfrac{1(1+1)}{2} = \frac{2}{2} = 1$ \\
    Тиймээс $n=1$ үед адилтгал үнэн байна.

    \item \textbf{Индукцийн таамаглал:} 
    $n=k$ үед адилтгал үнэн гэж үзье. Өөрөөр хэлбэл:
    \begin{equation}
    1 + 2 + 3 + \dots + k = \frac{k(k+1)}{2}
    \end{equation}

    \item \textbf{Индукцийн шилжилт ($n=k+1$):} 
    $n=k+1$ үед адилтгал үнэн болохыг буюу 
    $1 + 2 + \dots + k + (k+1) = \dfrac{(k+1)(k+2)}{2}$ болохыг баталъя.
    
    Индукцийн таамаглалыг ашиглавал:
    \begin{align*}
    1 + 2 + 3 + \dots + k + (k+1) &= (1 + 2 + \dots + k) + (k+1) \\
    &= \frac{k(k+1)}{2} + (k+1) \\
    &= \frac{k(k+1) + 2(k+1)}{2} \\
    &= \frac{(k+1)(k+2)}{2}
    \end{align*}
\end{enumerate}

$n=k+1$ үед үнэн болох нь батлагдлаа. Иймд математик индукцийн зарчмаар дурын натурал $n$ тооны хувьд адилтгал үнэн байна.
\end{prof}
\end{example}
\end{enumerate}
\newpage
\section*{Семинар 1}
\begin{enumerate}
  \item Дараах өгүүлбэрүүдийн аль нь хэллэг вэ? Хэллэг бол үнэнийг тогтоо: \\
  "7 нь тэгш"; $5<9$; "Хаалгаа хаа!"; $2+3=6$.
  \item $P:$ $"3<5"$, $Q:$ "7 нь тэгш" хэллэгүүд бол $P\land Q$, $P\lor Q$, $\lnot P$, $P\Rightarrow Q$, $Q\Rightarrow P$ хэллэгүүдийн үнэн худалыг тодорхойл.
  \item Үнэний хүснэгтээр шалга: $(P\Rightarrow Q)\Leftrightarrow(\lnot P\lor Q)$.
  \item Де Морганы хууль: $\lnot(P\land Q)\Leftrightarrow(\lnot P\lor\lnot Q)$-ийг хүснэгтээр батал.
  \item $F= P\land\lnot P$ хэллэг зөрчил байхыг батал.

  \item Батал: $P\land(P\lor Q)\Leftrightarrow P$ (absorption).
  \item $(P\Rightarrow Q)\lor(Q\Rightarrow P)$ нь тавтологи мөн эсэхийг тогтоо.
  \item $P:$ ``$n$ тэгш", $Q:$ ``$n^2$ 4-т хуваагдана". $P\Rightarrow Q$ хэллэгийг $n=6$ , $n=7$ ба $n=10$ тохиолдолд шалга. $P\Leftrightarrow Q$ теоремийг батал.
  \item $x>2 \Rightarrow x^2>4$ хэллэгийг монгол хэлээр бич.
  \item Үнэн хүснэгтээр шалга: $(P\land(Q\lor R))\Leftrightarrow((P\land Q)\lor(P\land R))$ (distributive).
  \item " Хэрэв тоо 6-д хуваагдвал 3-д хуваагдана" хэллэгийн урвуу хэллэгийг бичээд үнэн эсэхийг дүгнэ.
  \item Шууд батал: хэрэв $n$ сондгой бол $n^2$ сондгой.
  \item Эсрэгээс нь батал: хэрэв $ab$ сондгой бол $a$ ба $b$ сондгой.
  \item $a_1, a_2, \ldots a_n, \ldots$ дараалалын хувьд $\exists M > 0,  \forall n \in \mathbb{N}, (|a_n| < M)$ чанар биелдэг байг. Ийм дараалалын жишээ гарга. $\exists M > 0,  \forall n \in \mathbb{N}, (|a_n| < M)$ чанарын үгүйсгэл бич.
  \item Эквивалент: $(P\Rightarrow Q)\Leftrightarrow(\lnot Q\Rightarrow\lnot P)$.
  \item Шууд батал: хэрэв $a$ бүхэл бол $a^2+a$ тэгш.
  \item Дараах өгүүлбэрийн үгүйсгэлийг кванторыг зөв байршуулж бүрэн бич:\\
 $\forall y \in \mathbb{R}, \exists x \in \mathbb{R}, (f(x) = y)$.
  \item $\sqrt{6}$ иррационал болохыг батал.
  \item $4k+3$ хэлбэрийн анхны тоо хязгааргүй олон байдаг.
        (Заавар: эсрэгээс нь $p_1,\dots,p_m$ гэж аваад $N=4(p_1\cdots p_m)-1$-ийг авч үз.)
        \item $1^2+2^2+\cdots+n^2=\dfrac{n(n+1)(2n+1)}{6}$ батал. 
  \item $a,b$ нь дурын тоо, $n\in\mathbb{N}$ байг. 
\[
(a+b)^n=\sum_{k=0}^{n}\binom{n}{k}\,a^{\,n-k}b^{\,k},
\qquad
\binom{n}{k}=\frac{n!}{k!(n-k)!}\ (0\le k\le n).
\] индукцээр батал. 
  \item $a,b$ нь бодит тоо , $n\in\mathbb{N}$ байг. Дараах тэнцэтгэлийг бичээд батал:
\[
a^n-b^n=(a-b)(\sum_{k=0}^{n-1} a^{\,n-1-k} b^{\,k}).
\]
Мөн $n$ сондгой үед
\[
a^n+b^n=(a+b)(\sum_{k=0}^{n-1}(-1)^k a^{\,n-1-k} b^{\,k})
\]
байдгийг батал.
\item $a>b>0$ ба $n\in\mathbb{N}$ байг. 
\[
a^n-b^n \ge\; n\,b^{\,n-1}(a-b).
\]  батал.
\end{enumerate}

\newpage
\section*{Гэрийн даалгавар 1 }
\begin{enumerate}
  \item Үнэний хүснэгтээр шалга: $(P\land Q)\Rightarrow P$.
  \item Үнэний хүснэгтээр шалга: $(P\lor Q)\Leftrightarrow(\lnot P\Rightarrow Q)$.
  \item Эсрэг жишээ ол: "Бүх бүхэл тоо анхны эсвэл тэгш тоо".
  \item Эсрэгээс нь батал: хэрэв $n^2$ сондгой бол $n$ сондгой.
  \item Эквивалент: $(P\Rightarrow Q)\Leftrightarrow(\lnot Q\Rightarrow\lnot P)$.
  \item Эсрэг жишээ: "Бүх тэгш тоо анхны тоо".
  \item Шууд батал: хэрэв $a,b$ тэгш бол $a+b$ тэгш.
  \item Эсрэг жишээ: "Хэрэв $ab$ тэгш бол $a$ ба $b$ хоёулаа тэгш".
  \item Батал: хэрэв $a$ 4-т хуваагддаг бол $a^2$ 8-д хуваагдана.
  \item Эсрэг жишээ гарга: "Бүх анхны тоо сондгой".
  \item Тавтологи болохыг батал: $((P\Rightarrow Q)\land(Q\Rightarrow R))\Rightarrow(P\Rightarrow R)$.
  \item Тавтологи болохыг батал (resolution): $\big((P\lor Q)\land(\lnot P\lor R)\land(\lnot Q\lor R)\big)\Rightarrow R$.
  \item Бүх анхны $p$-ийн хувьд, хэрэв $p\mid n^2$ бол $p\mid n$ болохыг батал.
  \item $\sqrt{2}+\sqrt{3}$ иррационал болохыг батал.
        (Заавар: эсрэгээс нь $\sqrt{2}+\sqrt{3}=\tfrac{a}{b}$ гэж үз.)
  \item Анхны тоо \emph{хязгааргүй олон} байдгийг батал.
  \item $a_1, a_2, \ldots a_n, \ldots$ дараалалын хувьд $\forall M > 0, \exists n \in \mathbb{N}, (a_n > M)$ чанар биелдэг байг. Ийм дараалалын жишээ гарга.
   \item Дараах өгүүлбэрийн үгүйсгэлийг кванторыг зөв байршуулж бүрэн бич:\\
        $\forall\,\epsilon>0\, , \exists\,\delta>0\, , \big(|x-1|<\delta \Rightarrow |f(x)-f(1)|<\epsilon\big)$.
   
\end{enumerate}